\chapter{Esercizi: Interfacce e generici in Java}
\label{chap:esercizi-interfacce-generici-java}

\begin{exercise}
\label{ex:fi-esercizi-java-interfacce-design}
Qual è la differenza tra information hiding e modularità?
\end{exercise}

L'\emph{information hiding} nasconde i dettagli interni di un componente; la modularità organizza il programma in unità funzionali (\cref{sec:fi-java-interfacce-progettazione}).

\begin{exercise}
\label{ex:fi-esercizi-java-interfacce-object}
Quali sono gli scopi di \texttt{toString()} e \texttt{equals()} di \texttt{Object}?
\end{exercise}

\texttt{toString()} fornisce una rappresentazione testuale dell'oggetto; \texttt{equals()} ne verifica l'uguaglianza, che nell'implementazione di \texttt{Object} coincide con l'identità (\cref{sec:fi-java-interfacce-object}).

\begin{exercise}
\label{ex:fi-esercizi-java-interfacce-tostring}
Perché può essere opportuno sovrascrivere \texttt{toString()}?
\end{exercise}

Per ottenere una descrizione comprensibile delle informazioni significative, anziché la rappresentazione di base (\cref{sec:fi-java-interfacce-tostring}).

\begin{exercise}
\label{ex:fi-esercizi-java-interfacce-stampa}
Quando si utilizza implicitamente \texttt{toString()}?
\end{exercise}

In particolare nella stampa di un oggetto tramite \texttt{println(Object)} e nella concatenazione con una stringa (\cref{sec:fi-java-interfacce-tostring}).

\begin{exercise}
\label{ex:fi-esercizi-java-interfacce-equals}
Qual è la differenza tra \texttt{==} ed \texttt{equals()}?
\end{exercise}

Applicato a due riferimenti, \texttt{==} verifica se indicano lo stesso oggetto; \texttt{equals()} applica il criterio di uguaglianza definito dalla classe (\cref{sec:fi-java-interfacce-equals}).

\begin{exercise}
\label{ex:fi-esercizi-java-interfacce-class}
Che cosa restituiscono \texttt{getClass()} e \texttt{getName()}?
\end{exercise}

\texttt{getClass()} restituisce un oggetto \texttt{Class} che rappresenta la classe effettiva dell'istanza; \texttt{getName()}, invocato su tale oggetto, ne restituisce il nome completo, comprensivo dell'eventuale package (\cref{sec:fi-java-interfacce-tostring}).

\begin{exercise}
\label{ex:fi-esercizi-java-interfacce-accesso}
Quali sono gli ambiti di visibilità di \texttt{private}, \texttt{protected}, package e \texttt{public}?
\end{exercise}

Rispettivamente la classe dichiarante; le classi dello stesso package e le sottoclassi secondo le regole di Java; le classi dello stesso package; ogni classe che possa accedere alla classe dichiarante (\cref{tab:fi-java-interfacce-accesso}).

\begin{exercise}
\label{ex:fi-esercizi-java-interfacce-protected}
Perché un campo \texttt{protected} può compromettere l'incapsulamento?
\end{exercise}

Le sottoclassi e le classi del package possono dipendere direttamente dalla rappresentazione interna (\cref{sec:fi-java-interfacce-accesso}).

\begin{exercise}
\label{ex:fi-esercizi-java-interfacce-final}
Che cosa impedisce \texttt{final} applicato a metodi e classi?
\end{exercise}

La sovrascrittura del metodo e, rispettivamente, l'estensione della classe (\cref{tab:fi-java-interfacce-final}).

\clearpage

\begin{exercise}
\label{ex:fi-esercizi-java-interfacce-interfaccia}
Che cosa rappresenta un'interfaccia Java e come differisce da una classe?
\end{exercise}

Specifica un contratto comune senza costruttori o stato di istanza; una classe può contenere stato e consentire la costruzione di oggetti (\cref{def:fi-java-interfaccia}).

\begin{exercise}
\label{ex:fi-esercizi-java-interfacce-implements}
Come si implementa un'interfaccia?
\end{exercise}

Si dichiara \texttt{implements}; una classe concreta deve fornire i metodi astratti richiesti non già implementati (\cref{sec:fi-java-interfacce-implementazione}).

\begin{exercise}
\label{ex:fi-esercizi-java-interfacce-tipi-metodi}
Qual è la differenza tra metodi astratti, \texttt{static} e \texttt{default}?
\end{exercise}

I primi non hanno corpo; quelli \texttt{static} appartengono all'interfaccia; i \texttt{default} offrono una realizzazione di istanza ereditabile (\cref{sec:fi-java-interfacce-implementazione}).

\begin{exercise}
\label{ex:fi-esercizi-java-interfacce-variabile}
Si può dichiarare una variabile di tipo interfaccia e costruirne direttamente un oggetto?
\end{exercise}

La variabile è ammessa e può indicare un oggetto di una classe che implementa l'interfaccia, ma non è valido \texttt{new NomeInterfaccia()} (\cref{sec:fi-java-interfacce-riferimenti}).

\begin{exercise}
\label{ex:fi-esercizi-java-interfacce-metodi-disponibili}
Quali metodi sono utilizzabili attraverso un riferimento a un'interfaccia?
\end{exercise}

I metodi di istanza accessibili dichiarati o ereditati dall'interfaccia e i metodi pubblici di \texttt{Object}, ma non i metodi specifici della classe concreta non previsti dal tipo dichiarato (\cref{sec:fi-java-interfacce-riferimenti}).

\begin{exercise}
\label{ex:fi-esercizi-java-interfacce-polimorfismo}
Come si manifesta il polimorfismo con \texttt{Comparable}?
\end{exercise}

Il tipo statico abilita \texttt{compareTo()}, mentre il metodo eseguito dipende dall'oggetto effettivamente referenziato (\cref{sec:fi-java-interfacce-riferimenti,sec:fi-java-interfacce-comparable}).

\begin{exercise}
\label{ex:fi-esercizi-java-interfacce-risultato-compareto}
Quali valori può restituire \texttt{compareTo()}?
\end{exercise}

Un valore negativo se precede, zero se equivalente nell'ordine e positivo se segue; non necessariamente \(-1\) e \(1\) (\cref{def:fi-java-interfacce-comparable}).

\begin{exercise}
\label{ex:fi-esercizi-java-interfacce-generici}
Che cosa si intende per programmazione generica?
\end{exercise}

La definizione di classi e metodi che usano parametri di tipo specializzati dal contesto o dall'utilizzo (\cref{def:fi-java-interfacce-generici}).

\begin{exercise}
\label{ex:fi-esercizi-java-interfacce-vantaggi}
Quali vantaggi hanno i generici rispetto a \texttt{Object} e ai cast?
\end{exercise}

Permettono di riusare codice con verifica statica dei tipi e riducono i cast espliciti al recupero dei dati (\cref{sec:fi-java-interfacce-generici}).

\begin{exercise}
\label{ex:fi-esercizi-java-interfacce-sintassi}
Come si dichiara una classe generica con parametro \texttt{T}?
\end{exercise}

Con \lstinline|class Box<T> { ... }|, dove \texttt{T} è utilizzabile nei tipi dei campi, dei parametri e dei risultati (\cref{sec:fi-java-interfacce-generici}).

\begin{exercise}
\label{ex:fi-esercizi-java-interfacce-lettere}
Che cosa indicano per convenzione \texttt{T}, \texttt{E}, \texttt{K}, \texttt{V} e \texttt{N}?
\end{exercise}

Tipo, elemento, chiave, valore e tipo numerico. Sono convenzioni e non identificatori riservati (\cref{sec:fi-java-interfacce-generici}).

\begin{exercise}
\label{ex:fi-esercizi-java-interfacce-box-param}
Che differenza c'è tra \lstinline|Box<T>| e \lstinline|Box<String>|?
\end{exercise}

Il primo usa il parametro di tipo della dichiarazione; il secondo specifica un tipo concreto in un utilizzo (\cref{sec:fi-java-interfacce-generici}).

\begin{exercise}
\label{ex:fi-esercizi-java-interfacce-box-string}
Perché \lstinline|Box<String>| non può ricevere un \texttt{Integer}?
\end{exercise}

Il metodo \texttt{set()} richiede un \texttt{String}, e il controllo del tipo avviene in compilazione (\cref{sec:fi-java-interfacce-generici}).

\begin{exercise}
\label{ex:fi-esercizi-java-interfacce-metodo-generico}
Come si dichiara un metodo generico?
\end{exercise}

Si dichiara il parametro di tipo \lstinline|<T>| prima del tipo di ritorno, per esempio in \lstinline|static <T> void printArray(T[] elementi)| (\cref{sec:fi-java-interfacce-metodi-generici}).

\begin{exercise}
\label{ex:fi-esercizi-java-interfacce-array-algs}
Come si adatta \texttt{selectionSort()} a oggetti confrontabili?
\end{exercise}

Si utilizzano elementi confrontabili tramite \texttt{Comparable}; il confronto numerico è sostituito dal controllo del segno restituito da \texttt{compareTo()}, senza modificare la logica dell'algoritmo (\cref{sec:fi-java-interfacce-selection}).
